%% ============================================================
\section{Introduction}\label{sec:intro}
%% ============================================================

Loop invariants summarize loop behavior and reduce unbounded executions to
finite proof obligations, enabling verification beyond the executions covered
by finite testing.  Their value extends beyond proving a given
postcondition: legacy code can lack explicit assertions or specifications,
yet useful invariants can capture properties for reuse in later analysis and
verification.  When other loops or function calls follow a loop, the
function's final postcondition does not determine what should hold
immediately after that loop.  Proving the postcondition then requires an
intermediate assertion at that point, and a strong invariant, combined with
the loop's exit conditions, supplies exactly such a summary.  These settings
motivate discovering
informative loop summaries from program semantics without requiring a given
verification target.

Existing learning-based approaches, including Code2Inv, CLN2Inv, and
LIPuS~\citep{si2020code2inv,ryan2020cln2inv,yu2023lipus}, and LLM-based
methods such as LaM4Inv, Clause2Inv, and
LORIS~\citep{wu2024lam4inv,cao2025clause2inv,li2026loris}, use given
verification targets or target-derived feedback to guide invariant search.
We instead study \emph{target-hidden loop verification}:
generation and training scores use the loop context alone, while target
assertions and postconditions \(\mathcal G\) are restored only for final
verification.  Two challenges arise.  First, without a target, proof success
cannot directly guide candidate selection or training.  Inductiveness alone
is insufficient: \(\mathsf{true}\) is a valid invariant but supplies no useful
information about loop behavior.  Ideally, an invariant should include all
reachable states while excluding as many unreachable states as possible.
However, the exact reachable set is generally difficult to obtain, making
this precision difficult to measure directly as a training reward.
We therefore need a computable, target-independent proxy for invariant
strength.

Second, an LLM response may contain both useful
and incorrect clauses, while complementary constraints may be scattered
across responses.  Rejecting or rewarding each response as a whole obscures
the value of its useful parts.  A clause's value also depends on its companions:
two clauses that are not inductive individually can form an inductive
conjunction, with each helping establish that the other remains true after
an iteration.  Checking clauses independently can therefore discard useful
combinations.  Generation and training must account for their joint value.

We introduce \sam{} to align compositional inference, supervision labels,
and RL rewards.  At inference, it decomposes responses into clauses, pools
them for joint verification, repeatedly removes unproved clauses, and
conjoins the survivors.  During training, we construct candidate negative
traces by perturbing observed states and continuing execution beyond observed
loop exits.  With correctness checked by induction, rejecting more candidate
traces rules out more sampled behaviors.  Their rejection rate therefore
provides a computable strength proxy, not an exact measure of reachable-set
precision.
SFT uses verified multi-response
compositions that meet a coverage threshold as single-response training
labels; RL rewards the retained clauses' coverage and their contributions
to coverage across the response group.  Useful parts can thus receive
credit even when the complete response fails.  The two stages aim to make
multi-response discoveries accessible with fewer samples.  All supervision
comes from open-weight model generations, program execution, and
verification: SFT labels are synthesized once from Qwen3-8B rollouts and
shared across backbones, RL rewards are computed from each trained model's
own rollouts, and no closed model, human labeler, or external reward model
is used.

On 832 C programs, we restore hidden targets only for final verification to
test whether the generated invariants contain sufficient information to
complete a proof.  Compositional inference substantially improves base-model
verification, and aligned training yields further gains.  Our primary metric
is compose@1, the primary verification rate after filtering and composing one
response. We use pass@1 only as a standalone-response baseline.

%% Figure 1 retains the five-stage pipeline and all three training panels.
\begin{figure}[t]
\centering
% Shared vector artwork for Figure 1 and its standalone preview.
% Coordinate units follow the manuscript width; text is never scaled.
\begingroup
\definecolor{ovInk}{HTML}{25392F}
\definecolor{ovMuted}{HTML}{52635B}
\definecolor{ovBorder}{HTML}{D4DDD8}
\definecolor{ovGreen}{HTML}{2E7D5B}
\definecolor{ovRed}{HTML}{B85C47}
\definecolor{ovPaper}{HTML}{F7F9F8}
\pgfmathsetlengthmacro{\ovunit}{\textwidth/16}
% Helvetica is scaled by 0.92 in the manuscript: these yield 8/8.5 pt text.
\def\ovbody{\sffamily\fontsize{8.7}{10.2}\selectfont}
\def\ovtitle{\sffamily\fontsize{9.2}{10.8}\selectfont\bfseries}
\begin{tikzpicture}[
  x=\ovunit, y=.94*\ovunit, font=\ovbody, text=ovInk,
  >=Stealth,
  flow/.style={-{Stealth[length=1.35mm,width=.9mm]},
    draw=ovMuted, line width=.65pt},
  feedback/.style={-{Stealth[length=1.35mm,width=.9mm]},
    draw=ovGreen!85!black, line width=.65pt,
    dash pattern=on 2pt off 1.4pt, rounded corners=1.5pt},
  credit/.style={-{Stealth[length=.95mm,width=.65mm]},
    draw=ovMuted!72, line width=.45pt, dash pattern=on 1.4pt off 1pt},
  panel/.style={draw=ovBorder, fill=white, rounded corners=2pt,
    line width=.6pt},
  card/.style={draw=ovBorder, fill=white, rounded corners=1.5pt,
    line width=.5pt, align=center, inner xsep=3pt, inner ysep=4pt,
    outer sep=0pt},
  heading/.style={font=\ovtitle, text=ovInk, align=center},
  note/.style={font=\ovbody, text=ovMuted, align=center},
  chip/.style={rounded corners=.6pt, inner sep=0pt,
    minimum width=.30*\ovunit, minimum height=.30*\ovunit,
    line width=.4pt, font=\ovbody},
  a/.style={chip, fill=ovGreen!88, draw=ovGreen!88, text=white},
  b/.style={chip, fill=ovGreen!20, draw=ovGreen!40, text=ovInk},
  c/.style={a}, d/.style={b}, e/.style={a}, f/.style={b},
  u/.style={chip, fill=ovRed!22, draw=ovRed!55, text=ovRed!85!black},
  v/.style={u},
  pool/.style={rounded corners=1.5pt, draw=ovMuted!65, fill=white,
    line width=.55pt, dash pattern=on 2.3pt off 1.5pt},
  state/.style={circle, draw=ovGreen, fill=ovGreen,
    minimum size=.18*\ovunit, inner sep=0pt},
  candidate/.style={circle, draw=ovRed, fill=ovRed!16,
    line width=.65pt, minimum size=.20*\ovunit, inner sep=0pt}
]

% (a) Preserve the five-stage decompose--compose design.
\begin{scope}[shift={(0,4)}]
\draw[panel] (.10,3.72) rectangle (15.90,7.05);
\node[heading, anchor=west] at (.38,6.75)
  {(a)\; Decompose \(\rightleftarrows\) compose};
\draw[flow] (7.00,6.78)--(7.45,6.78);
\node[note, anchor=west] at (7.58,6.78) {data flow};
\draw[feedback] (9.65,6.78)--(10.10,6.78);
\node[note, anchor=west] at (10.23,6.78) {score / credit};
\draw[ovBorder, line width=.4pt] (.40,6.48)--(15.60,6.48);
\node[u] at (14.06,6.78) {\(u\)};
\node[note, anchor=west] at (14.31,6.78) {rejected};

\node[card, text width=2.30*\ovunit] (masked) at (1.45,5.12)
  {Masked loop\\[2pt]\(\mathcal L\! =\! (\mathrm{Init},B,S)\)\\[2pt]
   \textcolor{ovRed}{\(\mathcal G\) hidden}};
\draw[flow] (masked.east)--(2.77,5.12);

\fill[ovPaper, rounded corners=2pt] (2.86,3.96) rectangle (8.21,6.32);
\node[note, anchor=west] at (3.04,6.12) {within a rollout};
\node[heading] at (4.31,5.75) {1\; Rollout};
\node[heading] at (6.96,5.75) {2\; Decompose};
\foreach \y/\lab/\clauses in {
  5.38/1/{a,b,c,a,b},
  4.94/2/{c,u,d,e,c},
  4.50/3/{v,d,u,e,f}}{
  \node[card, minimum width=2.38*\ovunit,
    minimum height=.39*\ovunit, inner sep=0pt] at (4.32,\y) {};
  \node[note] at (3.40,\y) {\(y_{\lab}\)};
  \foreach \letter [count=\i from 0] in \clauses {
    \node[\letter] at (3.78+\i*.33,\y) {\(\letter\)};
    \node[\letter] at (6.55+\i*.34,\y) {\(\letter\)};
  }
  \node[note, anchor=east] at (6.24,\y) {\(A_{\lab}\)};
}
\draw[ovGreen, line width=.8pt]
  (5.29,5.38)--(5.34,5.32)--(5.44,5.47);
\foreach \y in {4.94,4.50}{
  \draw[ovRed, line width=.8pt]
    (5.29,\y-.07)--(5.43,\y+.07)
    (5.29,\y+.07)--(5.43,\y-.07);
}
\draw[flow] (5.56,4.94)--(5.75,4.94);
\draw[flow] (8.25,4.94)--(8.58,4.94);

\fill[ovPaper, rounded corners=2pt] (8.66,3.96) rectangle (15.60,6.32);
\node[note, anchor=west] at (8.87,6.12) {across rollouts};
\node[heading] at (9.68,5.75) {3\; Pool};
\node[heading] at (11.32,5.75) {4\; Filter};
\node[heading] at (13.29,5.75) {5\; Compose};
\node[pool, minimum width=1.53*\ovunit,
  minimum height=1.17*\ovunit] at (9.68,4.94) {};
\foreach \dx/\dy/\letter in {
  -.44/.39/a,0/.39/b,.44/.39/u,
  -.44/0/c,0/0/d,.44/0/v,
  -.22/-.39/e,.22/-.39/f}{
  \node[\letter] at (9.68+\dx,4.94+\dy) {\(\letter\)};
}
\draw[flow] (10.50,4.94)--(11.17,4.94);
\draw[ovMuted, line width=1.2pt] (11.23,4.56)--(11.23,5.32);
\draw[ovMuted!45, line width=.65pt] (11.35,4.56)--(11.35,5.32);
\draw[flow] (11.45,4.94)--(12.31,4.94);
\foreach \dx/\dy/\letter in {
  -.38/.21/a,0/.21/b,.38/.21/c,
  -.38/-.21/d,0/-.21/e,.38/-.21/f}{
  \node[\letter] at (12.98+\dx,4.94+\dy) {\(\letter\)};
}
\node[card, draw=ovGreen!40, fill=ovGreen!5,
  text width=1.25*\ovunit] (survivors) at (14.64,4.94)
  {Output\\[1pt]\(C^{\mathrm{surv}}\)};
\draw[flow] (13.57,4.94)--(survivors.west);
\end{scope}

% Shared output and original SFT/RL branches.
\draw[ovMuted, line width=.65pt]
  (survivors.south)--(14.64,7.39)--(2.12,7.39);
\foreach \x in {2.12,13.70}{
  \fill[ovMuted] (\x,7.39) circle (1.2pt);
  \draw[flow] (\x,7.39)--(\x,7.08);
}
\node[note, anchor=west] at (2.29,7.19) {composed set};
\node[note, anchor=west] at (13.87,7.19) {rollout group};

% (b) SFT: original acceptance, resampling, targets, and fine-tuning.
\draw[panel] (.10,.10) rectangle (3.96,7.05);
\node[heading, anchor=west] at (.38,6.76) {(b)\; SFT};
\node[note] at (2.03,6.17) {compose \(\to\) target};
\draw[ovBorder, line width=.4pt] (.38,6.48)--(3.68,6.48);
\node[card, text width=2.82*\ovunit] (composed) at (2.12,5.14)
  {Composed set\\[1pt]\(C_{\mathcal L}^{(t)}\)};
\node[card, text width=2.62*\ovunit] (accept) at (2.12,3.49)
  {Coverage \(\ge\tau\)?};
\draw[flow] (composed.south)--(accept.north);
\draw[flow, rounded corners=2pt]
  (accept.west)--(.43,3.49)--(.43,5.14)--(composed.west);
\node[note, anchor=west] at (.55,4.20) {resample};
\node[card, fill=ovGreen, draw=ovGreen, text=white,
  font=\ovtitle, inner xsep=6pt] (sfttarget) at (2.12,2.02) {SFT targets};
\draw[flow] (accept.south)--(sfttarget.north);
\node[note, anchor=west] at (2.24,2.77) {yes};
\node[card, text width=2.82*\ovunit] (sftmodel) at (2.12,.74) {Fine-tune};
\draw[flow] (sfttarget.south)--(sftmodel.north);

% (c) Shared scorer. Binary is a reference, not a sampler input.
\draw[panel] (4.30,.10) rectangle (11.16,7.05);
\node[heading, anchor=west] at (4.58,6.76) {(c)\; Target-free scorer};
\draw[ovBorder, line width=.4pt] (4.58,6.48)--(10.88,6.48);
\node[note] at (7.73,6.12)
  {\textbf{Binary}\quad valid set \(\mapsto\{0,1\}\)};

% A small, separate inset shows the two candidate-negative mechanisms.
\fill[ovPaper, rounded corners=2pt] (4.52,3.90) rectangle (10.93,5.90);
\node[heading, anchor=west] at (4.68,5.63) {Candidate negatives};
\node[note, anchor=west] at (4.68,4.95) {off-trajectory};
\node[note, anchor=west] at (4.68,4.16) {post-exit};
\foreach \x [count=\i] in {7.09,7.66,8.23}{
  \node[state] (off\i) at (\x,4.86) {};
}
\draw[flow, draw=ovGreen] (off1)--(off2);
\draw[flow, draw=ovGreen] (off2)--(off3);
\node[candidate] (offneg) at (8.24,5.31) {};
\draw[flow, draw=ovRed, dashed] (off2)--(offneg);
\node[state] (post1) at (7.09,4.12) {};
\node[state] (post2) at (7.66,4.12) {};
\draw[flow, draw=ovGreen] (post1)--(post2);
\draw[ovMuted, line width=.55pt, dashed] (8.03,3.98)--(8.03,4.37);
\node[note] at (8.03,4.50) {exit};
\node[candidate] (post3) at (8.41,4.12) {};
\node[candidate] (post4) at (9.01,4.12) {};
\draw[flow, draw=ovRed] (post2)--(post3);
\draw[flow, draw=ovRed] (post3)--(post4);
\node[circle, draw=ovRed!65, fill=white,
  minimum size=.49*\ovunit, inner sep=0pt] (negset) at (10.29,4.79)
  {\(\mathcal N\)};
\draw[flow, draw=ovRed!80, rounded corners=2pt]
  (offneg.east)--(9.35,5.31)--(negset.north west);
\draw[flow, draw=ovRed!80] (post4.east)--(negset.south west);
\draw[flow, rounded corners=2pt]
  (negset.south)--(10.29,3.77)--(8.03,3.77)--(8.03,3.66);

% Coverage retained at the same level as the SFT acceptance gate.
\node[heading, anchor=west] (coverageheading) at (4.58,3.49) {Coverage};
\draw[ovMuted!65, fill=white, line width=.55pt]
  (7.02,3.33) rectangle (9.02,3.64);
\fill[ovGreen!60] (7.02,3.33) rectangle (8.28,3.64);
\node[note, anchor=west] at (9.21,3.49) {\(\operatorname{cov}(V_i)\)};
\draw[feedback] (coverageheading.west)--(accept.east);
\draw[flow] (8.03,3.23)--(8.03,3.06);

% Original Shapley matrix, with more space for shared and unique coverage.
\node[heading, anchor=west, align=left] at (4.58,2.40) {Shapley\\allocation};
\foreach \y/\lab/\marks in {
  2.89/1/{1,1,1,0,0},
  2.47/2/{1,1,1,0,1},
  2.05/3/{0,0,1,2,1}}{
  \node[note, anchor=east] at (7.03,\y) {\(y_{\lab}\)};
  \foreach \mark [count=\i from 0] in \marks {
    \pgfmathsetmacro{\cx}{7.29+\i*.37}
    \ifcase\mark
      \node[circle, draw=ovBorder, fill=white, minimum size=.25*\ovunit,
        inner sep=0pt] at (\cx,\y) {};
    \or
      \node[circle, draw=ovGreen!65, fill=ovGreen!18,
        minimum size=.25*\ovunit, inner sep=0pt] at (\cx,\y) {};
      \draw[ovGreen, line width=.45pt]
        (\cx-.075,\y-.075)--(\cx+.075,\y+.075)
        (\cx-.075,\y+.075)--(\cx+.075,\y-.075);
    \or
      \node[circle, draw=ovGreen, fill=ovGreen!85,
        minimum size=.25*\ovunit, inner sep=0pt] at (\cx,\y) {};
    \fi
  }
  \node[note, anchor=west] at (9.19,\y) {\(\phi_{\lab}\)};
}
\node[circle, draw=ovGreen, fill=ovGreen!85,
  minimum size=.23*\ovunit, inner sep=0pt] at (7.08,1.62) {};
\node[note, anchor=west] at (7.27,1.62) {unique};
\node[circle, draw=ovGreen!65, fill=ovGreen!18,
  minimum size=.23*\ovunit, inner sep=0pt] at (8.73,1.62) {};
\draw[ovGreen, line width=.45pt]
  (8.66,1.55)--(8.80,1.69) (8.66,1.69)--(8.80,1.55);
\node[note, anchor=west] at (8.92,1.62) {shared};
\node[card, draw=ovGreen!40, fill=ovGreen!5,
  text width=5.54*\ovunit, inner ysep=2.5pt] (rewardformula) at (7.73,.74)
  {\textbf{Full reward per rollout}\\[1pt]
   \(r_i=\operatorname{cov}(V_i)+0.3\phi_i\)};

% (d) Preserve pooled survivor ownership and the many-to-many credit graph.
\draw[panel] (11.50,.10) rectangle (15.90,7.05);
\node[heading, anchor=west] at (11.78,6.76) {(d)\; RL};
\node[note] at (13.70,6.17) {attribute exploration};
\draw[ovBorder, line width=.4pt] (11.78,6.48)--(15.62,6.48);
\foreach \y/\lab/\clauses in {
  5.73/1/{a,b,c,a,b},
  5.27/2/{c,u,d,e,c},
  4.81/3/{v,d,u,e,f}}{
  \node[note, anchor=east] at (12.48,\y) {\(y_{\lab}\)};
  \foreach \letter [count=\i from 0] in \clauses {
    \node[\letter] at (12.77+\i*.41,\y) {\(\letter\)};
  }
}
\draw[flow] (13.60,4.59)--(13.60,4.40);
\draw[pool] (12.24,3.61) rectangle (15.07,4.39);
\foreach \letter [count=\i from 0] in {a,b,c,d,e,f}{
  \node[\letter] (owner\letter) at (12.54+\i*.44,3.84) {\(\letter\)};
}
\node[note] at (13.67,4.20) {pool, filtered once};
\node[note, text=ovInk] (r1) at (12.71,2.70) {\(r_1\)};
\node[note, text=ovInk] (r2) at (13.70,2.70) {\(r_2\)};
\node[note, text=ovInk] (r3) at (14.69,2.70) {\(r_3\)};
\foreach \letter/\target in {a/1,b/1,c/1,c/2,d/2,d/3,e/2,e/3,f/3}{
  \draw[credit] (owner\letter.south)--(r\target.north);
}
\draw[ovMuted, line width=.6pt]
  (r1.south)--(12.71,2.39)--(14.69,2.39)--(r3.south);
\draw[ovMuted, line width=.6pt] (r2.south)--(13.70,2.39);
\node[card, text width=3.65*\ovunit, inner ysep=3pt] (grpo) at (13.70,1.18)
  {\textbf{GRPO}\\[1pt]group-normalized\\advantage\\[1pt]
   \(\to\) policy update};
\draw[flow] (13.70,2.39)--(grpo.north);
\draw[feedback]
  (rewardformula.east)--(11.33,.74)--(11.33,2.70)--(r1.west);

\end{tikzpicture}%
\endgroup

\caption{\textbf{Decompose--compose is the organizing principle.}
(a) Pooling and joint filtering retain useful clauses from failed responses;
letters track clause identity. (b) SFT distills coverage-accepted compositions.
(c) Execution-derived candidate negatives supply coverage; Shapley allocation
splits shared rejection credit. (d) Retained clauses credit their proposing
rollouts for GRPO. The sampler inset shows off-trajectory perturbations and
post-exit body continuations. The hidden target \(\mathcal G\) is restored
only for final verification.}
\label{fig:overview}
\end{figure}


\begin{samepage}
\tightparagraph{Contributions.}
\begin{itemize}
  \item We introduce CRAFT, aligning target-hidden compositional inference
  with supervision and reward. A shared pooled filtering procedure recovers
  useful clauses from imperfect responses, constructs SFT labels, and
  identifies the contributions rewarded by RL. With one response, base
  Qwen3-8B reaches 37.86\% compose@1 versus 6.37\% pass@1.
  \item We turn multi-response discoveries into single-response learning
  signals. SFT learns accepted compositions, raising Qwen3-8B compose@1 to
  63.94\%; RL rewards retained clauses' coverage and group contributions,
  including those from failed responses, further raising it to 69.23\%.
  \item Our experiments reveal that whole-response success and compositional
  verification can diverge.  SFT broadens verification coverage and compresses
  the sampling budget; after substantial SFT, RL mainly improves sampling
  efficiency at small budgets, with smaller coverage gains at the largest
  tested budget.
\end{itemize}
\end{samepage}
